\clearpage
\section{Studium wykonalności}

    \subsection{Analiza trendów}
    
        W tym rodziale momencie zbiore dane z 27 dronów danej klasy następnie wykonam wykresy kluczowych danych,
        Wyznacze zależności najważniejszwgo parametru tej klasy dorana czyl czasu lotu od roku produkcji,
        Następnie wyznaczę wsznaczę zależności najważniejszych parametrów osiągowych od wydłużenia skrzydła,
        obciążenia powierzchni i ciągu (mocy) od tego parametru, Rownież postaram się użyć innych parametrów.
        Dane do wyresów załaczona w załaczniku Nr.1.

        \begin{figure}[!h]
            \centering \includegraphics[width=1\linewidth]{Wykresy/Czas_lotu_w_zależności_od_roku_produkcji.png}
            \caption{Wykres zależności czasu od roku produkcji}
        \end{figure}

        \begin{figure}[!h]
            \centering \includegraphics[width=1\linewidth]{Wykresy/Czas_lotu_w_zależności_od_obciążenia_mocy.png}
            \caption{Wykres zależności obciążenia mocy od czasu lotu}
        \end{figure}

        \begin{figure}[!h]
            \centering \includegraphics[width=1\linewidth]{Wykresy/Czas_lotu_w_zależności_od_obciążenia_powierzchni.png}
            \caption{Wykres zależności obciążenia powierzchni od czasu lotu}
        \end{figure}

        \begin{figure}[!h]
            \centering \includegraphics[width=1\linewidth]{Wykresy/Czas_lotu_w_zależności_od_wydłużenia.png}
            \caption{Wykres zależności obciążenia powierzchni od czasu lotu}
        \end{figure}

        \begin{figure}[!h]
            \centering \includegraphics[width=1\linewidth]{Wykresy/Masa_własna_startowa_w_zależności_od_masy_startowej.png}
            \caption{Wykres zależności wspołczunika masy własnej do startowej w zależnosći od masy startowej}
        \end{figure}
    
    \clearpage
    \subsubsection{Podsumowanie predykcji}
        Na podstawie danych z wykresów, Jestem wstanie oszacaować,
        minimalną wartość najważniejszego parametru czasu lotu,
        Również pozostałe współczuniki oreślahjace dynamikę BSP,
        Z wykresu Wykres zależności wspołczunika masy własnej do startowej w zależnosći od masy startowej jestem wstanie oszacować współczyniki A i C
        lepiej opisane w literaturze \cite{galinski}, 
        Współczyniki te nastepniesprawdziłe zwartościm podanumi w literaturz \cite{raymer}

        \begin{itemize}
            \item Czas lotu: 42,135 h
            \item Wydłużenie pałata: 15,468
            \item Obciązenie powierzchni: 97,177 kg/m\textsuperscript{2}
            \item Obciązenie mocy: 11,648 kg/kW
            \item Współczynik A: 1,510
            \item Współczynik C: -0,135
        \end{itemize}
    
    \subsection{Analiza konkurecjnych konstrukcji}
        Przed wykonaniem odrecznego rysunku prototypo zdecydowałem się na jeszcze jedenk krok,
        Mianowicie z badanej grupu dronów wyselekcjonowałem 5 najbardziej zblizonych konstrukcje pod zgledem parametrów opisanychw podzrodzielae podsumowanie predykcji,
        Selekcji dokonałem za pomoca odległosci euklidesowej opisanej w literaturze \cite{analnaInza}, 
        Funkcja dostepna na \cite{GE},
        Oraz zaprezentuje dodatkowedane nie  ujete wcześniej,

        \clearpage
        \subsubsection{Milkor 380}
            \begin{figure}[!h]
                \centering \includegraphics[width=1\linewidth]{Drony/Dron1/1.jpg}
                \caption{Milkor 380\cite{Dron1}}
            \end{figure}

        \begin{longtable}{| c | m{0,5\linewidth} | c |}
            \caption{Szczegółowe parametry Milkor 380}
            \label{table:milkor_380_specs} \\
            \hline
            \textbf{Lp} & \multicolumn{1}{c|}{\textbf{Nazwa parametru}} & \textbf{Wartość} \\ \hline \hline
            \endfirsthead
            
            \hline
            \endfoot
            
            1 & Producent & Milkor \\ \hline
            2 & Kraj produkcji & RPA \\ \hline
            3 & Dopasowanie & 93,23\% \\ \hline
            4 & Model silnika & Rotax 915 iS \\ \hline
            6 & Zasięg maksymalny  & 4000,0 km \\ \hline
            7 & Prędkość przelotowa / maksymalna & 150 / 250 km/h \\ \hline
            8 & Pułap operacyjny & 9144 m \\ \hline
            9 & Zasięg operacyjny SATA & 4000 km \\ \hline
            10 & Zasięg operacyjny LOS & 220 km \\ \hline
            11 & Długość / Wysokość / Rozpiętość & 9,00 / 2,40 / 18,60 m \\ \hline
            15 & Profil skrzydła root / tip & NACA 64-415 / NACA 64-412 \\ \hline
            16 & Masa własna & 700,0 kg \\ \hline
            17 & Masa startowa & 1300,0 kg \\ \hline
            18 & Maksymalna ładowność & 210,0 kg \\ \hline

        \end{longtable}

        \clearpage
        \subsubsection{Elbit Hermes 900}
            \begin{figure}[!h]
                \centering \includegraphics[width=1\linewidth]{Drony/Dron2/1.jpg}
                \caption{Elbit Hermes 900\cite{Dron2}}
            \end{figure}

            \begin{longtable}{| c | m{0,5\linewidth} | c |}
                \caption{Szczegółowe parametry Elbit Hermes 900}
                \label{table:hermes_900_specs} \\
                \hline
                \textbf{Lp} & \multicolumn{1}{c|}{\textbf{Nazwa parametru}} & \textbf{Wartość} \\ \hline \hline
                \endfirsthead
                
                \hline
                \endfoot
                
                1 & Producent & Elbit Systems \\ \hline
                2 & Kraj produkcji & Izrael \\ \hline
                3 & Dopasowanie & 89,99\% \\ \hline
                4 & Model silnika & Rotax 914 \\ \hline
                5 & Czas lotu & 36,0 h \\ \hline
                6 & Zasięg maksymalny & 5000,0 km \\ \hline
                7 & Prędkość przelotowa / maksymalna & 112 / 220 km/h \\ \hline
                8 & Pułap operacyjny & 9144 m \\ \hline
                9 & Zasięg operacyjny SATA & 5000 km \\ \hline
                10 & Zasięg operacyjny LOS & 250 km \\ \hline
                11 & Długość / Wysokość / Rozpiętość & 8,30 / 2,30 / 15,00 m \\ \hline
                12 & Wydłużenie płata & 15,00 \\ \hline
                13 & Obciążenie powierzchni & 78,67 kg/m$^2$ \\ \hline
                14 & Obciążenie mocy & 13,95 kg/kW \\ \hline
                15 & Profil skrzydła root / tip & Wortmann FX 63-137 \\ \hline
                16 & Masa własna & 750,0 kg \\ \hline
                17 & Masa startowa & 1180,0 kg \\ \hline
                18 & Maksymalna ładowność & 350,0 kg \\ \hline
            \end{longtable}

            \clearpage
            \subsubsection{Baykar Bayraktar TB2}

                \begin{figure}[!h]
                    \centering \includegraphics[width=1\linewidth]{Drony/Dron3/1.jpg}
                    \caption{Baykar Bayraktar TB2\cite{Dron3}}
                \end{figure}

                \begin{longtable}{| c | m{0,5\linewidth} | c |}
                    \caption{Szczegółowe parametry Baykar Bayraktar TB2}
                    \label{table:bayraktar_tb2_specs} \\
                    \hline
                    \textbf{Lp} & \multicolumn{1}{c|}{\textbf{Nazwa parametru}} & \textbf{Wartość} \\ \hline \hline
                    \endfirsthead
                    
                    \hline
                    \endfoot
                    
                    1 & Producent & Baykar Teknoloji \\ \hline
                    2 & Kraj produkcji & Turcja \\ \hline
                    3 & Dopasowanie & 89,35\% \\ \hline
                    4 & Model silnika & Rotax 912 \\ \hline
                    5 & Czas lotu & 27,0 h \\ \hline
                    6 & Zasięg maksymalny & 4000,0 km \\ \hline
                    7 & Prędkość przelotowa / maksymalna & 130 / 220 km/h \\ \hline
                    8 & Pułap operacyjny & 8229 m \\ \hline
                    9 & Zasięg operacyjny SATA & 4000 km \\ \hline
                    10 & Zasięg operacyjny LOS & 300 km \\ \hline
                    11 & Długość / Wysokość / Rozpiętość & 6,50 / 2,20 / 12,00 m \\ \hline
                    12 & Wydłużenie płata & 13,71 \\ \hline
                    13 & Obciążenie powierzchni & 66,67 kg/m$^2$ \\ \hline
                    14 & Obciążenie mocy & 9,52 kg/kW \\ \hline
                    15 & Profil skrzydła root / tip & NACA 4412 \\ \hline
                    16 & Masa własna & 420,0 kg \\ \hline
                    17 & Masa startowa & 700,0 kg \\ \hline
                    18 & Maksymalna ładowność & 150,0 kg \\ \hline
                \end{longtable}

            \clearpage
            \subsubsection{Shahed 129}

                \begin{figure}[!h]
                    \centering \includegraphics[width=1\linewidth]{Drony/Dron4/1.png}
                    \caption{Shahed 129\cite{Dron4}}
                \end{figure}

                \begin{longtable}{| c | m{0,5\linewidth} | c |}
                    \caption{SzczegółoAwe parametry Shahed 129}
                    \label{table:Shahed 129} \\
                    \hline
                    \textbf{Lp} & \multicolumn{1}{c|}{\textbf{Nazwa parametru}} & \textbf{Wartość} \\ \hline \hline
                    \endfirsthead
                    
                    \hline
                    \endfoot
                    1 & Producent & IAMIC \\ \hline
                    2 & Kraj produkcji & Iran \\ \hline
                    3 & Dopasowanie & 88,30\% \\ \hline
                    4 & Model silnika & Rotax 914 \\ \hline
                    5 & Czas lotu & 24,0 h \\ \hline
                    6 & Zasięg maksymalny & 1700,0 km \\ \hline
                    7 & Prędkość przelotowa / maksymalna & 150 / 175 km/h \\ \hline
                    8 & Pułap operacyjny & 7300 m \\ \hline
                    9 & Zasięg operacyjny SATA & 1700 km \\ \hline
                    10 & Zasięg operacyjny LOS & 200 km \\ \hline
                    11 & Długość / Wysokość / Rozpiętość & 8,00 / 3,10 / 16,00 m \\ \hline
                    12 & Wydłużenie płata & 14,88 \\ \hline
                    13 & Obciążenie powierzchni & 63,95 kg/m$^2$ \\ \hline
                    14 & Obciążenie mocy & 13,01 kg/kW \\ \hline
                    15 & Masa własna & 700,0 kg \\ \hline
                    16 & Masa startowa & 1100,0 kg \\ \hline
                    17 & Maksymalna ładowność & 400,0 kg \\ \hline
                \end{longtable}

            \clearpage
            \subsubsection{Wing Loong I}

                \begin{figure}[!h]
                    \centering \includegraphics[width=1\linewidth]{Drony/Dron5/1.jpg}
                    \caption{Shahed 129\cite{Dron5}}
                \end{figure}

                \begin{longtable}{| c | m{0,5\linewidth} | c |}
                    \caption{SzczegółoAwe parametry Shahed 129}
                    \label{table:Wing Loong I} \\
                    \hline
                    \textbf{Lp} & \multicolumn{1}{c|}{\textbf{Nazwa parametru}} & \textbf{Wartość} \\ \hline \hline
                    \endfirsthead
                    
                    \hline
                    \endfoot
                    1 & Producent & CAIG \\ \hline
                    2 & Kraj produkcji & Chiny \\ \hline
                    3 & Dopasowanie & 87,77\% \\ \hline
                    4 & Model silnika & Rotax 914 \\ \hline
                    5 & Czas lotu & 20,0 h \\ \hline
                    6 & Zasięg maksymalny & 4000,0 km \\ \hline
                    7 & Prędkość przelotowa / maksymalna & 150 / 280 km/h \\ \hline
                    8 & Pułap operacyjny & 7500 m \\ \hline
                    9 & Zasięg operacyjny SATA & 4000 km \\ \hline
                    10 & Zasięg operacyjny LOS & 200 km \\ \hline
                    11 & Długość / Wysokość / Rozpiętość & 9,05 / 2,70 / 14,00 m \\ \hline
                    12 & Wydłużenie płata & 12,17 \\ \hline
                    13 & Obciążenie powierzchni & 68,32 kg/m$^2$ \\ \hline
                    14 & Obciążenie mocy & 13,01 kg/kW \\ \hline
                    15 & Profil skrzydła root / tip & Drela GW-19 / Drela GW-27 \\ \hline
                    16 & Masa własna & 600,0 kg \\ \hline
                    17 & Masa startowa & 1100,0 kg \\ \hline
                    18 & Maksymalna ładowność & 200,0 kg \\ \hline
                \end{longtable}

    \clearpage
    \subsection{Model prototypu}

        \begin{figure}[!h]
            \centering \includegraphics[width=1\linewidth]{Model/model.pdf}
            \caption{Model prototypu}
        \end{figure}

        W strukturze projektowanego BSP zdecydowano się na zastosowanie układu centropłata, 
        Wybór ten podyktowany jest przede wszystkim pozwoli osiągnąć optymalne osiągi aerodynamiczne oraz wysokią manewrowości
        oraz wyposazony w chowane podwozie,  
        Konfiguracja ta wiąże się z ograniczeniem przestrzeni użytecznej wewnątrz kadłuba,
        Omawiany dron nie jest przeznaczony do przewozu wielkogabarytowego ładunku płatnego, 
        lecz do obserwacji, 
        Wyposażenie stanowi jedynie zintegrowana głowica optoelektroniczna,
        
        \paragraph{}
        Niemniej jednak, mając na uwadze konieczność optymalizacji czasu trwania procesu projektowego oraz dążenie do minimalizacji ryzyka błędu, 
        zdecydowano się na implementację usterzenia klasycznego, 
        
        \paragraph{}
        W projekcie zastosowano układ ze śmigłem pchającym z tyłu kadłuba, 
        Decyzja ta wynika z konieczności zapewnienia niezakłóconego pola widzenia dla głowicy optoelektronicznej, 
        Przeniesienie napędu na ogon uwalnia przestrzeń na dziobie drona, 
        eliminując ryzyko wchodzenia łopat śmigła w kadr i powstawania zakłóceń obrazu, 
        Dodatkowo chroni to delikatną optykę przed zanieczyszczeniami i uszkodzeniami podczas lądowania,

    \clearpage
    \subsection{Definicji misji i założenia najważniejszych osiągów samolotu}

        \begin{figure}[!h]
            \centering \includegraphics[width=1\linewidth]{ProfilMisji/misja.png}
            \caption{Profil misji BSP}
        \end{figure}

        \begin{itemize}
            \item Start
            \item Wznoszenie do wyskosści przelotowej
                \begin{itemize}
                    \item prędkosci wznoszenia przelotowego: 28 m/s 
                    \item kat wznoszenia: 15\textdegree
                \end{itemize}
            \item Dolot
                \begin{itemize}
                    \item prędkosci przelotowa: 42 m/s 
                    \item wysokość prtzelotowa: 3000 m
                \end{itemize}
            \item Wznoszenie do wyskosści patrolowej 
                \begin{itemize}
                    \item prędkosci wznoszenia przelotowego: 28 m/s 
                    \item kat wznoszenia: 10\textdegree
                \end{itemize}
            \item Patrolowanie
                \begin{itemize}
                    \item prędkosci patrolowa: 21 m/s
                    \item czas patrolu: 80000 s
                    \item wysokość patrolowa zalezna od misji wykrycie/rozpoznanie/identyfikacja: 13000/5200/2600 m
                \end{itemize}
            \item Opadanie
                \begin{itemize}
                    \item prędkosci wznoszenia przelotowego: 28 m/s 
                    \item kat opadania: 5\textdegree
                \end{itemize}
            \item Powrót
                \begin{itemize}
                    \item prędkosci przelotowa: 42 m/s 
                    \item wysokość prtzelotowa: 3000 m
                \end{itemize}
            \item Oczekiwanie: 
                \begin{itemize}
                    \item prędkosci patrolowa: 21 m/s
                    \item czas patrolu: 900 s
                    \item wysokość patrolowa zalezna od misji wykrycie/rozpoznanie/identyfikacja: 13000/5200/2600 m
                \end{itemize}
            \item Opadanie:
                \begin{itemize}
                    \item prędkosci wznoszenia przelotowego: 21 m/s 
                    \item kat opadania: 2\textdegree
                \end{itemize}
            \item Ladowanie

        \end{itemize}

    \clearpage
    \subsection{Szacunkowa biegunowa analityczna}
        
        W tym rodziale obliczę współczynnik oporu samolotu przy 0 współczynnik oporu samolotu,
        Wykonam to za poprzez przekształcenie wzoru współczynnik oporu samolotu,
        Odczytanie niezbędnych wartości z wcześniejszej analizy trendu i odczytachy wartości z tabel,
        Wszystkie niezbędene tabel i oraz wzory mozna znalez w \cite{galinski}, 
        Wzor po przekształceniu zamieszczam poniżej:

        \begin{align}
            C_{x0} = \frac{\pi \cdot AR \cdot e}{4 \cdot K_{\max}^2}
        \end{align}

        Dla moich wczesniej wspomnianych parametrów wartiość biegunowej wyszła: 0,0098

    \subsection{Wyznaczenia masy startowej i masy własnej samolotu}
        Z racji iz projektnie słozy  do obserwacji i nie rzuca ładubku w trakcie lotu,
        Do obliczeń masy ładunku samolotu wybrałem methode nr 2 opisana w \cite{galinski},
        Założeni i informacjęi najprawdopowdopniej do projektu, masa zastosowanej głowicy GOD-1 firmy PCO S,A zakładam 70 kg.
        
        \subsubsection{Obliczenie stosunków mas przed i po każdej fazie lotu}
        \begin{itemize}
            \item Rozruch:
                \begin{align}
                    \frac{W_{1}}{W_{0}}=0,97
                \end{align}
            \item Rozbieg:
                \begin{align}
                    \frac{W_{2}}{W_{1}}=0,985
                \end{align}
            \item Wznoszenie:
                \begin{align}
                    \frac{W_{3}}{W_{2}}=0,999
                \end{align}           
            \item Dolot:
                \begin{align}
                    \frac{W_{4}}{W_{3}}=0,984
                \end{align} 
            \item Wznoszenie:
                \begin{align}
                    \frac{W_{5}}{W_{4}}=0,997
                \end{align}  
            \item Patrolowanie:
                \begin{align}
                    \frac{W_{6}}{W_{5}}=0,778
                \end{align}
            \item Opadanie:
                \begin{align}
                    \frac{W_{7}}{W_{6}}=0,993
                \end{align}  
            \item Powrót:
                \begin{align}
                    \frac{W_{8}}{W_{7}}=0.993
                \end{align}  
            \item Opadanie:
                \begin{align}
                    \frac{W_{9}}{W_{8}}=0.994
                \end{align}  
            \item Oczekiwanie:
                \begin{align}
                    \frac{W_{10}}{W_{9}}=0.999
                \end{align}  
            \item Opadanie
                \begin{align}
                    \frac{W_{11}}{W_{10}}=0.994
                \end{align}
            \item Ladowanie:
                \begin{align}
                    \frac{W_{12}}{W_{11}}= 0.995
                \end{align}
        \end{itemize}

        \begin{align}
            \frac{W_{f}}{W_{0}} = 0.308
        \end{align}

        \subsubsection{Podsumowanie obliczenia massy staretowej i pustego samolotu}

            Następnie na podstawie methody iteracjnej została wyznaczona nastepująca masa.
            masa obliczona 881 masa założona: 880.96 roznica około 1 promila wiec okk.

        

    










 
